\documentclass[a4paper,12pt]{article}

\usepackage[T2A]{fontenc}
\usepackage[utf8]{inputenc}
\usepackage[russian]{babel}
\usepackage{amsmath,amssymb,mathtools}
\usepackage{helvet}
\renewcommand{\familydefault}{\sfdefault}
\usepackage{icomma}
\usepackage{geometry}
\usepackage{microtype}
\usepackage{tikz}
\usetikzlibrary{calc}
\usepackage{tabularx,booktabs}
\usepackage{enumitem}
\usepackage{hyperref}
\usepackage{parskip}
\usepackage{fancyhdr}
\usepackage{setspace}
\usepackage{xcolor}

\geometry{margin=2.1cm}
\onehalfspacing

\definecolor{accent}{HTML}{008080}
\definecolor{darkaccent}{HTML}{004D4D}
\definecolor{textgray}{HTML}{333333}
\definecolor{softgray}{HTML}{6E7777}
\definecolor{lightaccent}{HTML}{EAF3F3}

\hypersetup{
    colorlinks=true,
    linkcolor=darkaccent,
    urlcolor=accent,
    pdfauthor={Лёвин Артём Александрович},
    pdftitle={Чек-лист: степени и степенные преобразования}
}

\setlength{\parindent}{0pt}
\setlist[itemize]{leftmargin=1.5em,itemsep=0.35em,topsep=0.45em}
\setlist[enumerate]{leftmargin=1.7em,itemsep=0.65em,topsep=0.5em}

\pagestyle{fancy}
\fancyhf{}
\fancyhead[L]{\small\textcolor{softgray}{Анастасия Павлова}}
\fancyhead[R]{\small\textcolor{softgray}{Степени}}
\fancyfoot[C]{\small\textcolor{softgray}{\thepage}}
\renewcommand{\headrulewidth}{0.4pt}
\renewcommand{\headrule}{\hbox to\headwidth{\color{accent!45}\leaders\hrule height \headrulewidth\hfill}}

\newcommand{\topic}[1]{%
    \section*{\textcolor{darkaccent}{#1}}%
    \addcontentsline{toc}{section}{#1}%
}

\newcommand{\subtopic}[1]{%
    \subsection*{\textcolor{accent}{#1}}%
}

\newcommand{\important}[1]{\textbf{\textcolor{darkaccent}{#1}}}

\begin{document}

% ============================================================
% ТИТУЛЬНЫЙ ЛИСТ
% ============================================================

\begin{titlepage}
\thispagestyle{empty}

\begin{center}

\vspace*{2.4cm}

{\Large\textcolor{accent}{\textbf{ЧЕК-ЛИСТ}}}

\vspace{0.7cm}

{\huge\bfseries
Степени и степенные\\[0.2cm]
преобразования
}

\vspace{0.7cm}

{\large
Основные свойства степеней, дробные и отрицательные показатели,\\
корни, разложение чисел и алгоритмы преобразования выражений
}

\vspace{1.7cm}

{\large\textcolor{darkaccent}{
\textbf{Подготовка к ЕГЭ по математике}
}}

\vspace{2.1cm}

{\large
Лёвин Артём Александрович\\[0.15cm]
\textbf{эксклюзивно для Анастасии Павловой}
}

\vspace{1.3cm}

{\large \today}

\vfill

\begin{tikzpicture}[x=1cm,y=1cm]
    \draw[
        darkaccent!45,
        line width=1.2pt
    ]
    (0,0)
    .. controls (2.2,0.65) and (4.0,-0.45) .. (6.1,0.05)
    .. controls (8.1,0.55) and (10.2,-0.25) .. (12.6,0.35);
\end{tikzpicture}

\vspace{0.5cm}

\end{center}
\end{titlepage}

% ============================================================
% 1
% ============================================================

\topic{1. Что важно понимать о степени}

Степень
\[
a^n
\]
состоит из двух основных элементов:

\[
\underbrace{a}_{\text{основание}}
\qquad
\underbrace{n}_{\text{показатель степени}}.
\]

Например,
\[
5^3
\]
означает произведение трёх одинаковых множителей:
\[
5^3=5\cdot5\cdot5.
\]

Для натурального показателя
\[
a^n=
\underbrace{a\cdot a\cdot\ldots\cdot a}_{n\text{ множителей}}.
\]

Особые случаи:
\[
a^1=a,
\]
а при \(a\ne0\)
\[
a^0=1.
\]

Например,
\[
7^0=1,
\qquad
2^0=1,
\qquad
(-15)^0=1.
\]

\important{Важно.}
Выражение \(0^0\) в школьной алгебре не определяют.

\vspace{0.5cm}

При работе с дробными и действительными показателями далее предполагается,
что все выражения определены. Если основание степени должно быть положительным,
это условие необходимо учитывать.

% ============================================================
% 2
% ============================================================

\newpage
\topic{2. Главные свойства степеней}

\subtopic{2.1. Умножение степеней с одинаковым основанием}

\[
\boxed{a^m\cdot a^n=a^{m+n}}
\]

Показатели \important{складываются}.

Пример:
\[
2^3\cdot2^5=2^{3+5}=2^8.
\]

Ещё один пример:
\[
5^3\cdot5^{10}=5^{13}.
\]

Формула работает и справа налево:
\[
a^{m+n}=a^m\cdot a^n.
\]

Например,
\[
7^{x+y}=7^x\cdot7^y.
\]

\important{Условия применения:}

\begin{itemize}
    \item основания степеней одинаковы;
    \item между степенями стоит умножение.
\end{itemize}

Так делать нельзя:
\[
a^m+a^n\ne a^{m+n}.
\]

% ------------------------------------------------------------

\subtopic{2.2. Деление степеней с одинаковым основанием}

При \(a\ne0\)
\[
\boxed{\frac{a^m}{a^n}=a^{m-n}}.
\]

Показатели \important{вычитаются}.

Пример:
\[
\frac{5^{12}}{5^7}=5^{12-7}=5^5.
\]

Комбинируем два свойства:
\[
\frac{5^2\cdot5^{10}}{5^{12}}
=
5^{2+10-12}
=
5^0
=
1.
\]

\important{Главная идея:}
сложное выражение полезно разбивать на несколько простых преобразований.

% ============================================================
% 3
% ============================================================

\newpage
\topic{3. Одинаковые показатели степени}

\subtopic{3.1. Произведение}

Если показатели одинаковы,
\[
\boxed{a^n b^n=(ab)^n}.
\]

Например,
\[
5^2\cdot3^2=(5\cdot3)^2=15^2.
\]

То же свойство можно применять справа налево:
\[
(ab)^n=a^n b^n.
\]

Для трёх множителей:
\[
a^n b^n c^n=(abc)^n.
\]

Для целых показателей формула применяется непосредственно.
При дробных показателях необходимо учитывать область определения.

% ------------------------------------------------------------

\subtopic{3.2. Частное}

При \(b\ne0\)
\[
\boxed{\frac{a^n}{b^n}
=
\left(\frac ab\right)^n}.
\]

Например,
\[
\frac{3^2}{5^2}
=
\left(\frac35\right)^2.
\]

В обратную сторону:
\[
\left(\frac ab\right)^n
=
\frac{a^n}{b^n}.
\]

\subtopic{3.3. Что нельзя делать}

Степень в общем случае \important{не распределяется по сумме или разности}:
\[
(a+b)^n\ne a^n+b^n,
\]
\[
(a-b)^n\ne a^n-b^n.
\]

Например,
\[
(2+3)^2=25,
\]
а
\[
2^2+3^2=13.
\]

Следовательно,
\[
(2+3)^2\ne2^2+3^2.
\]

\important{Практический ориентир:}
если внутри конструкции появилась сумма или разность, свойства произведения
и частного степеней напрямую применять нельзя.

% ============================================================
% 4
% ============================================================

\newpage
\topic{4. Степень степени}

Одно из важнейших свойств:
\[
\boxed{(a^m)^n=a^{mn}}.
\]

Показатели \important{перемножаются}.

Например,
\[
(2^3)^4=2^{12}.
\]

Ещё пример:
\[
(5^2)^{0,32}=5^{2\cdot0,32}=5^{0,64}.
\]

Это свойство особенно полезно после разложения составного числа.

Например,
\[
25^{0,32}
=
(5^2)^{0,32}
=
5^{0,64}.
\]

Тогда
\[
5^{0,36}\cdot25^{0,32}
=
5^{0,36}\cdot5^{0,64}
=
5^1
=
5.
\]

\important{Скобки здесь принципиальны:}
\[
(5^2)^{0,32}.
\]

Они показывают, что во внешнюю степень возводится всё выражение \(5^2\).

% ============================================================
% 5
% ============================================================

\newpage
\topic{5. Отрицательный показатель степени}

Для \(a\ne0\)
\[
\boxed{a^{-n}=\frac1{a^n}}.
\]

И, наоборот,
\[
\boxed{\frac1{a^n}=a^{-n}}.
\]

Примеры:
\[
3^{-5}=\frac1{3^5},
\]
\[
2^{-1}=\frac12,
\]
\[
10^{-3}=\frac1{10^3}.
\]

Если показатель явно не указан,
\[
a=a^1.
\]

Поэтому
\[
\frac1a=a^{-1},
\qquad a\ne0.
\]

\subtopic{Числитель, отличный от единицы}

Например,
\[
\frac2x
=
2\cdot\frac1x
=
2x^{-1},
\qquad x\ne0.
\]

Аналогично
\[
\frac{7}{a^3}=7a^{-3},
\qquad a\ne0.
\]

% ============================================================
% 6
% ============================================================

\newpage
\topic{6. Переворачивание дроби и знак показателя}

При \(a\ne0\), \(b\ne0\)
\[
\boxed{
\left(\frac ab\right)^n
=
\left(\frac ba\right)^{-n}
}.
\]

Эквивалентная запись:
\[
\boxed{
\left(\frac ab\right)^{-n}
=
\left(\frac ba\right)^n
}.
\]

Пример:
\[
\left(\frac23\right)^{10}
=
\left(\frac32\right)^{-10}.
\]

Если показатель равен единице,
\[
\frac23
=
\left(\frac32\right)^{-1}.
\]

Ещё пример:
\[
\left(\frac25\right)^{-3}
=
\left(\frac52\right)^3
=
\frac{125}{8}.
\]

\important{Запомните смысл:}
при переворачивании основания дроби знак показателя меняется.

% ============================================================
% 7
% ============================================================

\newpage
\topic{7. Корни как степени}

Для положительного основания \(a\)
\[
\boxed{
\sqrt[k]{a^n}=a^{\frac nk}
},
\qquad
k\in\mathbb N,\quad k\ge2.
\]

Главный алгоритм:

\[
\text{внутренний показатель}
\quad\div\quad
\text{показатель корня}.
\]

Например,
\[
\sqrt[5]{a^3}=a^{3/5}.
\]

Для квадратного корня показатель \(2\) обычно не пишется:
\[
\sqrt a=\sqrt[2]a=a^{1/2},
\qquad a\ge0.
\]

Примеры:
\[
\sqrt2=2^{1/2},
\]
\[
\sqrt[3]{10}=10^{1/3},
\]
\[
\sqrt{10^3}=10^{3/2}.
\]

\important{Полезная привычка:}
показатели лучше сохранять в виде обыкновенных дробей:
\[
\frac32,\qquad
\frac53,\qquad
\frac7{12}.
\]

Так дальнейшие преобразования обычно становятся прозрачнее.

% ============================================================
% 8
% ============================================================

\newpage
\topic{8. Простые и составные числа}

\subtopic{Простые числа}

Натуральное число \(p>1\) называется \important{простым},
если имеет ровно два положительных делителя:
\[
1
\quad\text{и}\quad
p.
\]

Примеры:
\[
2,\ 3,\ 5,\ 7,\ 11,\ 13,\ldots
\]

Число \(1\) не является ни простым, ни составным.

\subtopic{Составные числа}

Натуральное число больше \(1\) называется \important{составным},
если имеет более двух положительных делителей.

Например,
\[
10
\]
имеет делители
\[
1,\ 2,\ 5,\ 10.
\]

Поэтому \(10\) --- составное число.

\subtopic{Разложение на простые множители}

Например,
\[
20=2\cdot2\cdot5=2^2\cdot5.
\]

Для числа \(48\):
\[
48
=
2\cdot2\cdot2\cdot2\cdot3
=
2^4\cdot3.
\]

Разложение особенно полезно в задачах со степенями, поскольку позволяет
получить одинаковые основания.

Например,
\[
8=2^3,\qquad
9=3^2,\qquad
16=2^4,
\]
\[
25=5^2,\qquad
27=3^3,\qquad
81=3^4.
\]

% ============================================================
% 9
% ============================================================

\newpage
\topic{9. Главный алгоритм преобразования выражений}

При решении задач со степенями удобно двигаться в следующем порядке.

\begin{enumerate}
    \item \important{Разложите составные числа} и постарайтесь получить удобные основания.
    \item \important{Переведите корни в степени} с дробным показателем.
    \item \important{При необходимости избавьтесь от дробей} с помощью отрицательных показателей.
    \item Посмотрите, какие основания или показатели стали одинаковыми.
    \item Примените подходящее свойство степени.
    \item Сократите и упростите показатели.
    \item Повторяйте преобразования, пока выражение не станет простым.
    \item Выполните окончательный численный расчёт.
\end{enumerate}

Схематично:

\begin{center}
\begin{tikzpicture}[
    node distance=1.35cm,
    every node/.style={font=\small},
    step/.style={
        draw=darkaccent!70,
        rounded corners=3pt,
        inner sep=7pt,
        text width=10.6cm,
        align=center
    },
    arrow/.style={->,>=stealth,line width=0.8pt,darkaccent!70}
]

\node[step] (a) {Составные числа \(\longrightarrow\) удобные основания};
\node[step,below of=a] (b) {Корни \(\longrightarrow\) дробные показатели};
\node[step,below of=b] (c) {Дроби \(\longrightarrow\) отрицательные показатели, если это удобно};
\node[step,below of=c] (d) {Ищем одинаковые основания или одинаковые показатели};
\node[step,below of=d] (e) {Применяем свойства степеней и упрощаем};

\draw[arrow] (a)--(b);
\draw[arrow] (b)--(c);
\draw[arrow] (c)--(d);
\draw[arrow] (d)--(e);

\end{tikzpicture}
\end{center}

% ============================================================
% 10
% ============================================================

\newpage
\topic{10. Разбор типовых преобразований}

\subtopic{Пример 1. Составное основание}

Вычислим:
\[
5^{0,36}\cdot25^{0,32}.
\]

Разложим
\[
25=5^2.
\]

Тогда
\[
5^{0,36}\cdot(5^2)^{0,32}
=
5^{0,36}\cdot5^{0,64}.
\]

Складываем показатели:
\[
5^{0,36+0,64}=5^1=5.
\]

\[
\boxed{5}
\]

% ------------------------------------------------------------

\subtopic{Пример 2. Отрицательный показатель}

Вычислим:
\[
3^{6,5}\cdot9^{-2,25}.
\]

Так как
\[
9=3^2,
\]
получаем
\[
3^{6,5}\cdot(3^2)^{-2,25}.
\]

Степень степени:
\[
(3^2)^{-2,25}=3^{-4,5}.
\]

Следовательно,
\[
3^{6,5}\cdot3^{-4,5}
=
3^{2}
=
9.
\]

\[
\boxed{9}
\]

% ------------------------------------------------------------

\subtopic{Пример 3. Разложение составного знаменателя}

Вычислим:
\[
\frac{2^{3,5}\cdot3^{5,5}}{6^{4,5}}.
\]

Разложим
\[
6=2\cdot3.
\]

Тогда
\[
6^{4,5}
=
(2\cdot3)^{4,5}
=
2^{4,5}\cdot3^{4,5}.
\]

Следовательно,
\[
\frac{2^{3,5}\cdot3^{5,5}}
     {2^{4,5}\cdot3^{4,5}}
=
2^{3,5-4,5}\cdot3^{5,5-4,5}.
\]

Получаем
\[
2^{-1}\cdot3
=
\frac12\cdot3
=
\frac32.
\]

\[
\boxed{\frac32}
\]

% ------------------------------------------------------------

\subtopic{Пример 4. Корень и дробные показатели}

Вычислим:
\[
\left(
\frac{2^{1/3}\cdot2^{1/4}}
     {\sqrt[12]{2}}
\right)^2.
\]

Переведём корень в степень:
\[
\sqrt[12]{2}=2^{1/12}.
\]

Тогда
\[
\left(
2^{1/3+1/4-1/12}
\right)^2.
\]

Приводим показатели к знаменателю \(12\):
\[
\frac13+\frac14-\frac1{12}
=
\frac4{12}+\frac3{12}-\frac1{12}
=
\frac6{12}
=
\frac12.
\]

Получаем
\[
\left(2^{1/2}\right)^2
=
2.
\]

\[
\boxed{2}
\]

% ============================================================
% 11
% ============================================================

\newpage
\topic{11. Скобки: где они особенно важны}

Скобки помогают однозначно показать структуру выражения.

\subtopic{Подстановка}

Если вместо \(25\) подставляется \(5^2\), то в выражении
\[
25^{0,32}
\]
получаем
\[
(5^2)^{0,32}.
\]

Именно скобки показывают, что показатель \(0,32\) относится ко всему выражению
\(5^2\).

\subtopic{Отрицательное основание}

Необходимо различать:
\[
(-2)^4=16
\]
и
\[
-2^4=-16.
\]

Во втором случае степень относится только к числу \(2\).

\subtopic{Дробное основание}

Правильно:
\[
\left(\frac23\right)^5.
\]

Если основание представляет собой дробь, сумму, разность или произведение,
скобки почти всегда делают запись безопаснее и понятнее.

% ============================================================
% 12
% ============================================================

\newpage
\topic{12. Типичные ошибки}

\begin{enumerate}
    \item \textbf{Складывать показатели при разных основаниях.}

    Неверно:
    \[
    2^3\cdot3^5=6^8.
    \]

    Здесь показатели различаются, поэтому такая формула неприменима.

    \item \textbf{Складывать показатели при сложении степеней.}

    Неверно:
    \[
    a^3+a^5=a^8.
    \]

    Свойство \(a^m a^n=a^{m+n}\) относится к произведению.

    \item \textbf{Забывать перемножить показатели.}

    \[
    (a^m)^n=a^{mn}.
    \]

    Например,
    \[
    (5^2)^3=5^6.
    \]

    \item \textbf{Терять знак отрицательного показателя.}

    \[
    a^{-n}=\frac1{a^n},
    \qquad a\ne0.
    \]

    \item \textbf{Переворачивать дробь без изменения знака показателя.}

    Правильно:
    \[
    \left(\frac23\right)^5
    =
    \left(\frac32\right)^{-5}.
    \]

    \item \textbf{Неверно переводить корень в степень.}

    Правильно:
    \[
    \sqrt[k]{a^n}=a^{n/k}.
    \]

    \item \textbf{Забывать ограничения.}

    Например,
    \[
    a^{-2}
    \]
    требует
    \[
    a\ne0.
    \]

    Квадратный корень
    \[
    \sqrt a
    \]
    в действительных числах требует
    \[
    a\ge0.
    \]

    \item \textbf{Использовать десятичные показатели там, где удобнее дроби.}

    Например,
    \[
    1,5=\frac32.
    \]

    В задачах со степенями запись
    \[
    a^{3/2}
    \]
    обычно удобнее, чем
    \[
    a^{1,5}.
    \]
\end{enumerate}

% ============================================================
% 13
% ============================================================

\newpage
\topic{13. Краткая памятка}

\renewcommand{\arraystretch}{1.45}

\begin{table}[h]
\centering
\caption{Основные свойства степеней}
\small
\begin{tabularx}{\linewidth}{@{}X X@{}}
\toprule
\textbf{Ситуация} & \textbf{Преобразование} \\
\midrule

Одинаковые основания, умножение
&
\(\displaystyle a^m a^n=a^{m+n}\)
\\

Одинаковые основания, деление
&
\(\displaystyle \frac{a^m}{a^n}=a^{m-n},\ a\ne0\)
\\

Одинаковые показатели, произведение
&
\(\displaystyle a^n b^n=(ab)^n\)
\\

Одинаковые показатели, частное
&
\(\displaystyle \frac{a^n}{b^n}
=\left(\frac ab\right)^n,\ b\ne0\)
\\

Степень степени
&
\(\displaystyle (a^m)^n=a^{mn}\)
\\

Отрицательный показатель
&
\(\displaystyle a^{-n}=\frac1{a^n},\ a\ne0\)
\\

Переворот дроби
&
\(\displaystyle
\left(\frac ab\right)^{-n}
=
\left(\frac ba\right)^n
\)
\\

Корень
&
\(\displaystyle
\sqrt[k]{a^n}=a^{n/k}
\)
\\

Нулевая степень
&
\(\displaystyle a^0=1,\ a\ne0\)
\\

\bottomrule
\end{tabularx}
\end{table}

\vspace{0.7cm}

\important{Главный вопрос после каждого шага:}

\begin{center}
\large
\textbf{«Стало ли выражение проще и какое свойство можно применить теперь?»}
\end{center}

% ============================================================
% ТРЕНИРОВОЧНЫЙ БЛОК
% ============================================================

\newpage
\topic{Тренировочный блок. Путь Б}

\textcolor{softgray}{
Выполняйте задания последовательно. Формулы разрешается держать перед глазами.
Главная цель --- выработать навык распознавания подходящего свойства.
}

\subtopic{27.08.2026}

\begin{enumerate}
    \item Вычислите:
    \[
    \frac{2^7\cdot2^5}{2^9}.
    \]

    \item Вычислите:
    \[
    \frac{5^8}{5^3}.
    \]

    \item Представьте одной степенью:
    \[
    3^4\cdot7^4.
    \]

    \item Представьте одной степенью:
    \[
    (2^3)^4.
    \]

    \item Вычислите:
    \[
    \frac{4^5\cdot2^3}{8^4}.
    \]
\end{enumerate}

% ------------------------------------------------------------

\newpage
\subtopic{28.08.2026}

\begin{enumerate}
    \item Представьте степенью с отрицательным показателем:
    \[
    \frac1{2^6}.
    \]

    \item Запишите в виде обыкновенной дроби:
    \[
    5^{-2}.
    \]

    \item Вычислите:
    \[
    \left(\frac37\right)^{-2}.
    \]

    \item Вычислите:
    \[
    \left(\frac25\right)^{-3}.
    \]

    \item Вычислите:
    \[
    2^{-3}\cdot4^2.
    \]
\end{enumerate}

% ------------------------------------------------------------

\newpage
\subtopic{29.08.2026}

\begin{enumerate}
    \item Представьте в виде степени:
    \[
    \sqrt7.
    \]

    \item Представьте в виде степени:
    \[
    \sqrt[3]{5^2}.
    \]

    \item Представьте в виде степени с основанием \(3\):
    \[
    \sqrt[4]{3^6}.
    \]

    \item Вычислите, предварительно представив подкоренное число степенью двойки:
    \[
    \sqrt{64}.
    \]

    \item Вычислите:
    \[
    \left(\sqrt[3]{2}\right)^6.
    \]
\end{enumerate}

% ------------------------------------------------------------

\newpage
\subtopic{30.08.2026}

\begin{enumerate}
    \item Разложите на простые множители и запишите компактно:
    \[
    72.
    \]

    \item Представьте в виде степени с основанием \(2\):
    \[
    \sqrt{32}.
    \]

    \item Представьте в виде степени с основанием \(3\):
    \[
    \frac1{81}.
    \]

    \item Вычислите:
    \[
    25^{3/2}.
    \]

    \item Вычислите:
    \[
    27^{-2/3}.
    \]
\end{enumerate}

% ------------------------------------------------------------

\newpage
\subtopic{31.08.2026 }

\begin{center}

\vspace{1.1cm}

{\LARGE\bfseries\textcolor{darkaccent}{
С Днём рождения, Анастасия!
}}

\vspace{10.8cm}

{\large
Желаю Вам ярких впечатлений,\\
интересных открытий и больших успехов в учёбе!
}

\vspace{1.3cm}

\begin{figure}[p]
    \centering
    \includegraphics[
        width=\textwidth,
        height=0.88\textheight,
        keepaspectratio
    ]{01.png}
\end{figure}

\vspace{1.2cm}

\textcolor{accent}{
\textbf{Сегодня --- без тренировочных заданий.}
}

\end{center}

% ------------------------------------------------------------

\newpage
\subtopic{01.09.2026}

\begin{enumerate}
    \item Вычислите:
    \[
    5^{0,36}\cdot25^{0,32}.
    \]

    \item Вычислите:
    \[
    3^{6,5}\cdot9^{-2,25}.
    \]

    \item Вычислите:
    \[
    \frac{2^{3,5}\cdot3^{5,5}}{6^{4,5}}.
    \]

    \item Вычислите:
    \[
    \left(
    \frac{2^{1/3}\cdot2^{1/4}}
         {2^{1/12}}
    \right)^2.
    \]

    \item Вычислите:
    \[
    16^{3/4}\cdot8^{-2/3}.
    \]
\end{enumerate}

% ------------------------------------------------------------

\newpage
\subtopic{02.09.2026}

\begin{enumerate}
    \item Вычислите:
    \[
    \frac{27^{2/3}\cdot9^{1/2}}{3^2}.
    \]

    \item Вычислите:
    \[
    \left(\frac1{32}\right)^{-2/5}.
    \]

    \item Вычислите:
    \[
    \left(\sqrt[4]{81}\cdot3^{-1}\right)^2.
    \]

    \item Вычислите:
    \[
    \frac{5^{7/3}\cdot25^{-2/3}}{5^{-1}}.
    \]

    \item Вычислите:
    \[
    \left(\frac49\right)^{-3/2}
    \cdot
    \left(\frac8{27}\right)^{2/3}.
    \]
\end{enumerate}

% ============================================================
% ОТВЕТЫ
% ============================================================

\newpage
\topic{Ответы к тренировочному блоку}

\renewcommand{\arraystretch}{1.55}

\begin{table}[h]
\centering
\caption{Ответы}
\small
\begin{tabularx}{\linewidth}{@{}l *{5}{>{\centering\arraybackslash}X}@{}}
\toprule
\textbf{Дата}
&
\textbf{1}
&
\textbf{2}
&
\textbf{3}
&
\textbf{4}
&
\textbf{5}
\\
\midrule

27.08
&
\(8\)
&
\(3125\)
&
\(21^4\)
&
\(2^{12}\)
&
\(2\)
\\

28.08
&
\(2^{-6}\)
&
\(\dfrac1{25}\)
&
\(\dfrac{49}{9}\)
&
\(\dfrac{125}{8}\)
&
\(2\)
\\

29.08
&
\(7^{1/2}\)
&
\(5^{2/3}\)
&
\(3^{3/2}\)
&
\(8\)
&
\(4\)
\\

30.08
&
\(2^3\cdot3^2\)
&
\(2^{5/2}\)
&
\(3^{-4}\)
&
\(125\)
&
\(\dfrac19\)
\\

01.09
&
\(5\)
&
\(9\)
&
\(\dfrac32\)
&
\(2\)
&
\(2\)
\\

02.09
&
\(3\)
&
\(4\)
&
\(1\)
&
\(25\)
&
\(\dfrac32\)
\\

\bottomrule
\end{tabularx}
\end{table}

\vspace{1cm}

\begin{center}
\textcolor{darkaccent}{
\textbf{Если ответ не совпал, вернитесь к последнему шагу,
на котором выражение ещё было очевидно понятным.}
}
\end{center}

\end{document}